% Part: second-order-logic
% Chapter: sol-and-sets
% Section: power-of-continuum

\documentclass[../../../include/open-logic-section]{subfiles}

\begin{document}

\olfileid{sol}{set}{pow}

\olsection{Kardinalitas Kontinuum}

\begin{explain}
Ing logika orde kapindho kita bisa nguwantifikasi
subhimpunan !!{domain}, nanging ora bisa nguwantifikasi
himpunan kang anggotane subhimpunan-subhimpunan !!{domain}.
Kanggo nindakake iki kanthi langsung, kita mbutuhake
logika \emph{orde katelu}. Umpamane, yen arep nyatakake
Teorema Cantor yen ora ana fungsi !!{injective} saka
himpunan pangkat sawijining himpunan menyang himpunan iku
dhewe, kita bisa nyoba ngrumusake mangkene:
``kanggo saben himpunan~$X$ lan saben himpunan~$P$,
yen $P$ iku himpunan pangkat saka~$X$, mula
ora $\cardle{P}{X}$''. Kanggo nyatakake yen $P$
iku himpunan pangkat saka~$X$, kita kudu ngrumusake
yen !!{element}s saka $P$ iku persis kabeh subhimpunan
saka~$X$, dadi kurang luwih
$\lforall[Y][(P(Y) \liff Y \subseteq X)]$.
Masalahe ana ing $P(Y)$: iki dudu !!a{formula}
logika orde kapindho, amarga mung suku kang bisa dadi
argumen variabel relasi siji-papan kaya~$P$.

Nanging kita bisa \emph{nyimulasi} kuantifikasi
marang himpunan kang anggotane himpunan, yen
domaine cukup gedhe. Gagasane nggunakake kasunyatan
yen relasi rong-papan~$R$ nggayutake !!{element}s
!!{domain} karo !!{element}s !!{domain}.
Yen ana~$R$ kaya mangkono, kita bisa nglumpukake
kabeh !!{element}s kang gegayutan karo
sawijining~$x$ miturut~$R$: $\Setabs{y \in
  \Domain{M}}{R(x,y)}$ iku himpunan kang
``dikodhe'' dening~$x$. Kosok baline, yen
$Z \subseteq \Pow{\Domain{M}}$ iku sawijining
kulawarga subhimpunan saka~$\Domain{M}$,
lan cacah !!{element}s ing~$\Domain{M}$
paling ora padha karo cacah himpunan ing~$Z$,
mula uga ana relasi~$R \subseteq \Domain{M}^2$
saengga saben $Y \in Z$ dikodhe dening
sawijining~$x$ nganggo~$R$.
\end{explain}

\begin{defn}
Yen $R \subseteq \Domain{M}^2$, mula
$x$ \emph{ngodhe miturut~$R$}
$\Setabs{y \in \Domain{M}}{R(x, y)}$.
\end{defn}

Yen !!a{element}~$x \in \Domain{M}$ ngodhe
miturut~$R$ sawijining himpunan
$Z \subseteq \Domain{M}$, mula himpunan
$Y \subseteq \Domain{M}$ ngodhe himpunan kang
anggotane himpunan, yaiku himpunan-himpunan kang
dikodhe dening !!{element}s saka~$Y$. Dadi
himpunan $Y$ bisa ngodhe miturut~$R$
$\Pow{X}$. Iki kedadeyan yen lan mung yen kanggo
saben $Z \subseteq X$, ana $x \in Y$ kang
ngodhe miturut~$R$ himpunan~$Z$, lan saben
$x \in Y$ ngodhe miturut~$R$ sawijining~$Z \subseteq X$.

\begin{prop}
!!{formula}
  \[
  \fn{Codes}(x, R, Z) \ident \lforall[y][(Z(y) \liff
    R(x, y))]
  \]
nyatakake yen $s(x)$ ngodhe miturut~$s(R)$
himpunan~$s(Z)$. !!{formula}
\begin{multline*}
  \fn{Pow}(Y, R, X) \ident \\
  \lforall[Z][(Z \subseteq X \lif \lexists[x][(Y(x) \land
    \fn{Codes}(x, R, Z))])] \land {} \\
  \lforall[x][(Y(x) \lif
    \lforall[Z][(\fn{Codes}(x, R, Z) \lif Z \subseteq X)])]
\end{multline*}
nyatakake yen $s(Y)$ ngodhe miturut~$s(R)$
himpunan pangkat saka~$s(X)$, yaiku
unsur-unsur saka $s(Y)$ ngodhe miturut~$s(R)$ persis
kabeh subhimpunan saka $s(X)$.
\end{prop}

\begin{explain}
Kanthi cara iki kita bisa nyatakake pratelan ngenani
himpunan pangkat kanthi nguwantifikasi kode
subhimpunan, dudu subhimpunane dhewe. Umpamane,
Teorema Cantor saiki bisa dinyatakake kanthi nyebut
yen ora ana fungsi !!{injective} saka himpunan
unsur kang dadi kode kanggo himpunan pangkat
saka~$X$ menyang~$X$ dhewe.
\end{explain}

\begin{prop}
Ukara
\begin{multline*}
  \lforall[X][\lforall[Y][\lforall[R][(\fn{Pow}(Y, R, X) \lif \\
      \lnot \lexists[u][(\lforall[x][\lforall[y][(\eq[u(x)][u(y)] \lif
            \eq[x][y])]] \land {}\\
        \lforall[x][(Y(x) \lif X(u(x)))])])]]]
\end{multline*}
iku valid.
\end{prop}

\begin{explain}
Himpunan pangkat saka himpunan !!a{denumerable}
iku !!{nonenumerable}, mula kardinalitase luwih
gedhe tinimbang kardinalitas sembarang himpunan
!!{denumerable} (yaiku $\aleph_0$). Ukuran
$\Pow{\Nat}$ diarani ``kardinalitas kontinuum'',
amarga padha ukurane karo titik-titik ing garis
wilangan nyata, yaiku~$\Real$. Yen !!{domain}
cukup gedhe kanggo ngodhe himpunan pangkat
saka himpunan !!a{denumerable}, kita bisa
nyatakake yen sawijining himpunan ukurane
kontinuum kanthi nyebut yen himpunan iku
padha cacahe karo sawijining himpunan~$Y$
kang ngodhe himpunan pangkat saka himpunan~$X$
kang ukurane~$\aleph_0$. (Yen domaine ora
cukup gedhe, yaiku ora ngemot subhimpunan
kang padha cacahe karo~$\Real$, mula uga
ora ana relasi kang ngodhe~$\Pow{X}$.)
\end{explain}

\begin{prop}
Yen $\cardle{\Real}{\Domain{M}}$, mula
!!{formula}
\begin{multline*}
\fn{Cont}(Y) \ident
\lexists[X][\lexists[R][((\fn{Aleph}_0(X) \land
      \fn{Pow}(Y, R, X)) \land \\
      \lforall[x][\lforall[y][((Y(x) \land {}
      Y(y) \land \lforall[z][R(x,z) \liff R(y,z)]) \lif \eq[x][y])]])]]
\end{multline*}
nyatakake yen $\cardeq{s(Y)}{\Real}$.
\end{prop}

\begin{proof}
  $\fn{Pow}(Y, R, X)$ nyatakake yen $s(Y)$ ngodhe
  miturut~$s(R)$ himpunan pangkat
  saka~$s(X)$, kang miturut $\fn{Aleph}_0(X)$
  bisa dienumerasi. Dadi $s(Y)$ paling ora
  padha gedhene karo kontinuum, sanadyan bisa
  luwih gedhe yen sawetara unsur $s(Y)$ ngodhe
  subhimpunan saka~$X$ kang padha.
  Kahanan iki dicegah dening konjungsi pungkasan,
  kang mbutuhake sesambungan antarane unsur
  $s(Y)$ lan subhimpunan saka $s(X)$ lumantar
  $s(R)$ iku !!{injective}.
\end{proof}

\begin{prop}
$\cardeq{\Domain{M}}{\Real}$ yen lan mung yen
\begin{multline*}
  \Sat{M}{\lexists[Y][(\lforall[x][Y(x)] \land \fn{Cont}(Y))]}.
\end{multline*}
\end{prop}

\begin{explain}
Hipotesis Kontinuum iku pratelan yen ukuran
kontinuum minangka kardinalitas !!{nonenumerable}
kapisan, yaiku yen $\Pow{\Nat}$
ukurane~$\aleph_1$.
\end{explain}

\begin{prop}
Hipotesis Kontinuum bener yen lan mung yen
\[\fn{CH} \ident
\lforall[X][(\fn{Aleph}_1(X) \liff \fn{Cont}(X))]\]
valid.
\end{prop}

Gatekna yen ora bener yen $\lnot \fn{CH}$
valid yen lan mung yen Hipotesis Kontinuum
salah. Ing !!a{enumerable} domain, ora ana
subhimpunan kang ukurane~$\aleph_1$ lan
uga ora ana subhimpunan kang ukurane
kontinuum, mula $\fn{CH}$ tansah bener
ing !!a{enumerable} domain. Nanging kita
bisa menehi ukara liya kang valid yen lan
mung yen Hipotesis Kontinuum salah:

\begin{prop}
Hipotesis Kontinuum salah yen lan mung yen
\[\fn{NCH} \ident
\lforall[X][(\fn{Cont}(X) \lif \lexists[Y][(Y \subseteq X \land \lnot
    \fn{Count}(Y) \land \lnot \cardeq{X}{Y})])]\]
valid.
\end{prop}

\end{document}
